\setchaptertoc
\chapter{Evaluation}\label{chp:evaluation}
\csummary{Is algorithm B better than algorithm C?} %This command can be used to display short exemplary questions in the contents overview

\begin{summary}
We did\ldots
\end{summary}

\section[RQs]{Research Questions}\label{sec:evaluation:rqs}
What are your central research questions?

\begin{rqs}
    \item \label{rq:concept} What is the concept of this thesis?
    \item \label{rq:implementation} How is the concept implemented?
    \item \label{rq:evaluation} How is the implementation evaluated?
    \begin{rqs}
        \item \label{rq:evaluation:correctness} How is the correctness evaluated? \begin{rqs}
                                                                                      \item Please do not nest the questions that deep!
        \end{rqs}
        \item \label{rq:evaluation:performance} How is the performance evaluated?
    \end{rqs}
\end{rqs}

We can now reference \ref{rq:concept} or with cleveref: \cref{rq:concept}.

We can also reference sections, e.g., \cref{sec:evaluation:rqs}

\section{Method}
How did you perform your experiments?

\section{Results}
What are the results of your experiments?

\section{Discussion}
What are the answers to your research questions?

\section[Threats]{Threats to Validity}
What threatens the validity of your findings?
Discuss at least internal and external validity, possibly more~\cite{WRH+12}.

\section{Summary}
Short summary of this section